\chapter{Limitations}

\section{Campaign State and Reproducibility}

The report uses benchmark run artifacts for reported values, but reproducibility
still has constraints:

\begin{itemize}
\item headline main and operations results are from commit-locked \(N=5\) runs on
one host;
\item latest report manifests capture commit and clean-tree provenance, but older
supporting artifacts have weaker embedded provenance;
\item some comparisons (for example LowMC baseline vs optimised and AES backend
variants) rely partly on supporting runs outside the latest main benchmark
campaign.
\end{itemize}

These limits do not remove the measured signal, but they reduce statistical and
forensic confidence relative to repeated runs across more hosts and full target
coverage under one lock policy.

Their practical impact is different for different claims. Large effect-size
claims, such as LowMC remaining far slower than AES, are less sensitive to host
noise. Small wall-clock deltas, such as CTR versus CTR+HMAC total-time
differences, are much more sensitive and are therefore interpreted mainly through
cycle counts and statistical caution.

\section{Single-zkVM Scope}

All headline measurements are from one proving stack (RISC Zero). This project
does not claim cross-zkVM generality.

The conclusions most at risk of non-transfer are relative performance rankings
that depend heavily on trace shape (for example ``LowMC is much slower than AES
in this stack''). A different zkVM with different memory behavior, instruction
cost profile, or proving backend could shift those ratios.

By contrast, the methodological conclusion is more robust: primitive-level
intuition alone is insufficient, and deployment decisions should be backed by
stack-specific measurement.

\section{Security Argument Depth}

The security treatment is an explicit game-based proof sketch grounded in
standard assumptions. A full formal proof with tight reductions and mechanized
verification is out of scope.

This means the argument should be interpreted as structured assurance rather than
machine-checked theorem-level certainty. It also does not prove full
chosen-ciphertext AEAD security, replay resistance, sender identity, or
side-channel resistance.

\section{Transport Authentication Scope}

The implementation authenticates ciphertext data inside the guest pipeline, but
does not add an external proof-transport signature layer. As a result, current
claims cover confidentiality and integrity of cryptographic payload handling
under verified execution, but not sender identity, session binding, or replay
guarantees in distributed deployments where proof artifacts transit untrusted
channels.

\section{Impact and Mitigation Summary}

Table~\ref{tab:limitations-impact-mitigation} summarizes how each limitation
affects the report claims and how it can be mitigated.

\begin{table}[H]
\centering
\scriptsize
\begin{tabular}{p{0.24\linewidth}p{0.35\linewidth}p{0.31\linewidth}}
\hline
Limitation & Impact on results & Mitigation \\
\hline
Single-host \(N=5\) campaign plus legacy single-trial support snapshots & Limits
confidence for small timing deltas and mixed-provenance comparisons & Replicate
headline runs on more hosts and extend the lock policy to legacy support targets \\
Weaker provenance in older supporting artifacts & Reduces forensic
reproducibility for historical snapshots if code changes after the report & Embed
git revision and dirty-tree status in all trial and aggregate artifacts \\
Single-zkVM scope & Prevents general claims about all zkVMs & Port the same
workload contract to another zkVM and compare ranking stability \\
Proof sketch rather than full reduction & The Chapter 10 game is explicit, but
the reductions are informal and not tight & Expand Chapter 10 into theorem-level
reductions with tightness discussion \\
No transport identity layer & Does not prevent sender impersonation or replay in
a network protocol & Add signatures, session identifiers, and replay controls at
deployment layer \\
\hline
\end{tabular}
\caption{Limitations, impact on claims, and mitigation path.}
\label{tab:limitations-impact-mitigation}
\end{table}
